\documentclass[10pt,letterpaper]{article}
\usepackage[utf8]{inputenc}
\usepackage[intlimits]{amsmath}
\usepackage{amsfonts}
\usepackage{amssymb}
\usepackage{ragged2e}
\usepackage[letterpaper, margin=0.5in]{geometry}
\usepackage{graphicx}
\usepackage{cancel}
\usepackage{mathtools}
\usepackage{tabularx}
\usepackage{arydshln}
\usepackage{tensor}
\usepackage{array}
\usepackage{xcolor}
\usepackage[boxed]{algorithm}
\usepackage[noend]{algpseudocode}
\usepackage{listings}
\usepackage{textcomp}
% \usepackage[pdf,tmpdir,singlefile]{graphviz}
\usepackage{mathrsfs}
\usepackage{bbm}
\usepackage{tikz}
\usepackage{tikz-cd}
\usepackage{enumitem}
\usepackage{arydshln}
\usepackage{relsize}
\usepackage{multicol}
\usepackage{scalerel}
\usepackage{upgreek}
\usepackage{ifthen}

\usetikzlibrary{bayesnet}
\usetikzlibrary{calc,positioning}
\setlist{noitemsep}

%%%%%%%%%%%%%%%%%%%%%%%%%%%%%
% Formatting commands
%%%%%%%%%%%%%%%%%%%%%%%%%%%%%
\newcommand{\n}{\hfill\break}
\newcommand{\nn}{\vspace{0.5\baselineskip}\n}
\newcommand{\up}{\vspace{-\baselineskip}}
\newcommand{\hangblock}[2]{\par\noindent\settowidth{\hangindent}{\textbf{#1: }}\textbf{#1: }\nolinebreak#2}
\newcommand{\lemma}[1]{\hangblock{Lemma}{#1}}
\newcommand{\defn}[1]{\hangblock{Defn}{#1}}
\newcommand{\thm}[1]{\hangblock{Thm}{#1}}
\newcommand{\cor}[1]{\hangblock{Cor}{#1}}
\newcommand{\prop}[1]{\hangblock{Prop}{#1}}
\newcommand{\ex}[1]{\hangblock{Ex}{#1}}
\newcommand{\exer}[1]{\hangblock{Exer}{#1}}
\newcommand{\fact}[1]{\hangblock{Fact}{#1}}
\newcommand{\remark}[1]{\hangblock{Remark}{#1}}
\newcommand{\proven}{\;$\square$\n}
\newcommand{\problem}[1]{\par\noindent{\nolinebreak#1}\n}
\newcommand{\problempart}[2]{\par\noindent\indent{}\settowidth{\hangindent}{\textbf{(#1)} \indent{}}\textbf{(#1) }\nolinebreak#2\n}
\newcommand{\ptxt}[1]{\textrm{\textnormal{#1}}}
\newcommand{\inlineeq}[1]{\centerline{$\displaystyle #1$}}
\newcommand{\pageline}{\up\par\noindent\rule{\textwidth}{0.1pt}}

%%%%%%%%%%%%%%%%%%%%%%%%%%%%%
% Math commands
%%%%%%%%%%%%%%%%%%%%%%%%%%%%%
% Set Theory
\newcommand{\card}[1]{\left|#1\right|}
\newcommand{\set}[1]{\left\{#1\right\}}
\newcommand{\setmid}{\;\middle|\;}
\newcommand{\ps}[1]{\mathcal{P}\left(#1\right)}
\newcommand{\pfinite}[1]{\mathcal{P}^{\ptxt{finite}}\left(#1\right)}
\newcommand{\naturals}{\mathbb{N}}
\newcommand{\N}{\naturals}
\newcommand{\integers}{\mathbb{Z}}
\newcommand{\Z}{\integers}
\newcommand{\rationals}{\mathbb{Q}}
\newcommand{\Q}{\rationals}
\newcommand{\reals}{\mathbb{R}}
\newcommand{\R}{\reals}
\newcommand{\complex}{\mathbb{C}}
\newcommand{\C}{\complex}
\newcommand{\halfPlane}{\mathbb{H}}
\let\HSym\H
\let\H\relax
\newcommand{\H}{\halfPlane}
\newcommand{\comp}{^{\complement}}
\DeclareMathOperator{\Hom}{Hom}
\newcommand{\Ind}{\mathbbm{1}}
\newcommand{\cut}{\setminus}
\DeclareMathOperator{\elem}{elem}

% Differentiable Manifolds
\newcommand{\RP}{\mathbb{RP}}
\newcommand{\CP}{\mathbb{RP}}
\newcommand{\osubset}{\overset{\mathclap{\scalebox{0.5}{\ptxt{open}}}}{\subset}}
\newcommand{\osubseteq}{\overset{\mathclap{\scalebox{0.5}{\ptxt{open}}}}{\subseteq}}
\newcommand{\osupset}{\overset{\mathclap{\scalebox{0.5}{\ptxt{open}}}}{\supset}}
\newcommand{\osupseteq}{\overset{\mathclap{\scalebox{0.5}{\ptxt{open}}}}{\supseteq}}
\newcommand{\pdat}[3]{\left.\pd{#1}{#2}\right|_{#3}}
\DeclareMathOperator{\so}{so}
\DeclareMathOperator{\codim}{codim}
\DeclareMathOperator{\Diff}{Diff}
\let\dSym\d
\let\d\relax
\newcommand{\d}{\partial}
\DeclareMathOperator{\gl}{gl}
\DeclareMathOperator{\Ad}{Ad}

% Graph Theory
\let\deg\relax
\DeclareMathOperator{\deg}{deg}
\newcommand{\degp}{\ptxt{deg}^{+}}
\newcommand{\degn}{\ptxt{deg}^{-}}
\newcommand{\precdot}{\mathrel{\ooalign{$\prec$\cr\hidewidth\hbox{$\cdot\mkern0.5mu$}\cr}}}
\newcommand{\succdot}{\mathrel{\ooalign{$\cdot\mkern0.5mu$\cr\hidewidth\hbox{$\succ$}\cr\phantom{$\succ$}}}}
\DeclareMathOperator{\cl}{cl}
\DeclareMathOperator{\affdim}{affdim}

% Probability
\newcommand{\parSymbol}{\P}
\newcommand{\Prob}{\mathbb{P}}
\renewcommand{\P}{\Prob}
\newcommand{\Avg}{\mathbb{E}}
\newcommand{\E}{\Avg}
\DeclareMathOperator{\Var}{Var}
\DeclareMathOperator{\cov}{cov}
\DeclareMathOperator{\Unif}{Unif}
\DeclareMathOperator{\Binom}{Binom}
\newcommand{\CI}{\mathrel{\text{\scalebox{1.07}{$\perp\mkern-10mu\perp$}}}}
\DeclareMathOperator{\Ber}{Ber}
\DeclareMathOperator{\Bin}{Bin}
\DeclareMathOperator{\Geom}{Geom}
\DeclareMathOperator{\Poisson}{Poisson}

% Standard Math
\newcommand{\inv}{^{-1}}
\newcommand{\abs}[1]{\left|#1\right|}
\newcommand{\ceil}[1]{\left\lceil{}#1\right\rceil{}}
\newcommand{\floor}[1]{\left\lfloor{}#1\right\rfloor{}}
\newcommand{\conj}[1]{\overline{#1}}
\newcommand{\of}{\circ}
\newcommand{\tri}{\triangle}
\newcommand{\inj}{\hookrightarrow}
\newcommand{\surj}{\twoheadrightarrow}
\newcommand{\ndiv}{\nmid}
\renewcommand{\epsilon}{\varepsilon}
\newcommand{\divides}{\mid}
\newcommand{\ndivides}{\nmid}
\DeclareMathOperator{\lcm}{lcm}
\DeclareMathOperator{\sgn}{sgn}
\newcommand{\map}[4]{\!\!\!\begin{array}[t]{rcl}#1 & \!\!\!\!\to & \!\!\!\!#2\\ {}#3 & \!\!\!\!\mapsto & \!\!\!\!#4\end{array}}
\newcommand{\bigsum}[2]{\smashoperator[lr]{\sum_{\scalebox{#1}{$#2$}}}}
\DeclareMathOperator{\gcf}{gcf}
\newcommand{\restr}[2]{\left.#1\right|_{#2}}

% Linear Algebra
\newcommand{\Id}{\textrm{\textnormal{Id}}}
\newcommand{\im}{\textrm{\textnormal{im}}}
\newcommand{\norm}[1]{\abs{\abs{#1}}}
\newcommand{\tpose}{^{T}\!}
\newcommand{\iprod}[1]{\left<#1\right>}
\newcommand{\giprod}{\iprod{\;\,,\;}}
\DeclareMathOperator{\tr}{tr}
\DeclareMathOperator{\trace}{tr}
\newcommand{\chgBasMat}[3]{\!\!\tensor*[_{#1}]{\left[#2\right]}{_{#3}}}
\newcommand{\vecBas}[2]{\tensor*[]{\left[#1\right]}{_{#2}}}
\DeclareMathOperator{\GL}{GL}
\DeclareMathOperator{\Mat}{Mat}
\DeclareMathOperator{\vspan}{span}
\DeclareMathOperator{\rank}{rank}
\newcommand{\V}[1]{\vec{#1}}
\DeclareMathOperator{\proj}{proj}
\DeclareMathOperator{\compProj}{comp}
\DeclareMathOperator{\row}{row}
\newcommand{\smallPMatrix}[1]{\paren{\begin{smallmatrix}#1\end{smallmatrix}}}
\newcommand{\smallBMatrix}[1]{\brack{\begin{smallmatrix}#1\end{smallmatrix}}}
\newcommand{\pmat}[1]{\begin{pmatrix}#1\end{pmatrix}}
\newcommand{\bmat}[1]{\begin{bmatrix}#1\end{bmatrix}}
\newcommand{\dual}{^{*}}
\newcommand{\pinv}{^{\dagger}}
\newcommand{\horizontalMatrixLine}{\ptxt{\rotatebox[origin=c]{-90}{$|$}}}
\DeclareMathOperator{\range}{range}
\DeclareMathOperator{\Symm}{Symm}
\DeclareMathOperator{\SU}{SU}
\DeclareMathOperator{\U}{U}

% Multilinear Algebra
\let\Lsym\L
\let\L\relax
\DeclareMathOperator{\L}{\mathscr{L}}
\DeclareMathOperator{\A}{\mathcal{A}}
\DeclareMathOperator{\Alt}{Alt}
\DeclareMathOperator{\Sym}{Sym}
\newcommand{\ot}{\otimes}
\newcommand{\ox}{\otimes}
\DeclareMathOperator{\asc}{asc}
\DeclareMathOperator{\asSet}{set}
\DeclareMathOperator{\sort}{sort}
\DeclareMathOperator{\ringA}{\mathring{A}}
\DeclareMathOperator{\Sh}{Sh}
\DeclareMathOperator{\Bil}{Bil}

% Topology
\newcommand{\closure}[1]{\overline{#1}}
\newcommand{\uball}{\mathcal{U}}
\DeclareMathOperator{\Int}{Int}
\DeclareMathOperator{\Ext}{Ext}
\DeclareMathOperator{\Bd}{Bd}
\DeclareMathOperator{\rInt}{rInt}
\DeclareMathOperator{\ch}{ch}
\DeclareMathOperator{\ah}{ah}
\newcommand{\LargerTau}{\mathlarger{\mathlarger{\mathlarger{\mathlarger{\tau}}}}}
\newcommand{\Tau}{\mathcal{T}}

% Analysis
\DeclareMathOperator{\Graph}{Graph}
\DeclareMathOperator{\epi}{epi}
\DeclareMathOperator{\hypo}{hypo}
\DeclareMathOperator{\supp}{supp}
\newcommand{\lint}[2]{\underset{#1}{\overset{#2}{{\color{black}\underline{{\color{white}\overline{{\color{black}\int}}\color{black}}}}}}}
\newcommand{\uint}[2]{\underset{#1}{\overset{#2}{{\color{white}\underline{{\color{black}\overline{{\color{black}\int}}\color{black}}}}}}}
\newcommand{\alignint}[2]{\underset{#1}{\overset{#2}{{\color{white}\underline{{\color{white}\overline{{\color{black}\int}}\color{black}}}}}}}
\newcommand{\extint}{\ptxt{ext}\int}
\newcommand{\extalignint}[2]{\ptxt{ext}\alignint{#1}{#2}}
\newcommand{\conv}{\ast}
\newcommand{\pd}[2]{\frac{\partial{}#1}{\partial{}#2}}
\newcommand{\del}{\nabla}
\DeclareMathOperator{\grad}{grad}
\DeclareMathOperator{\curl}{curl}
\let\div\relax
\DeclareMathOperator{\div}{div}
\DeclareMathOperator{\vol}{vol}

% Complex Analysis
\let\Re\relax
\DeclareMathOperator{\Re}{Re}
\let\Im\relax
\DeclareMathOperator{\Im}{Im}
\DeclareMathOperator{\Res}{Res}

% Abstract Algebra
\DeclareMathOperator{\ord}{ord}
\newcommand{\generated}[1]{\left<#1\right>}
\newcommand{\cycle}[1]{\smallPMatrix{#1}}
\newcommand{\id}{\ptxt{id}}
\newcommand{\iso}{\cong}
\DeclareMathOperator{\Aut}{Aut}
\DeclareMathOperator{\SL}{SL}
\DeclareMathOperator{\op}{op}
\newcommand{\isom}[4]{\!\!\!\begin{array}[t]{rcl}#1 & \!\!\!\!\overset{\sim}{\to} & \!\!\!\!#2\\ #3 & \!\!\!\!\mapsto & \!\!\!\!#4\end{array}}
\newcommand{\F}{\mathbb{F}}
\newcommand{\acts}{\;\reflectbox{\rotatebox[origin=c]{-90}{$\circlearrowright$}}\;}
\newcommand{\disjunion}{\mathrel{\text{\scalebox{1.07}{$\perp\mkern-10mu\perp$}}}}
\DeclareMathOperator{\SO}{SO}
\DeclareMathOperator{\stab}{stab}
\DeclareMathOperator{\nullity}{nullity}
\DeclareMathOperator{\Perm}{Perm}
\DeclareMathOperator{\nsubgp}{\vartriangleleft}
\DeclareMathOperator{\notnsubgp}{\ntriangleleft}
\newcommand{\presentation}[2]{\left<#1\;\middle|\;#2\right>}
\DeclareMathOperator{\Char}{char}
\DeclareMathOperator{\fchar}{char}
\DeclareMathOperator{\triv}{triv}
\DeclareMathOperator{\reg}{reg}
\DeclareMathOperator{\std}{std}
\DeclareMathOperator{\Func}{Func}
\DeclareMathOperator{\End}{End}

% Convex Optimization
\let\sectionSymbol\S
\let\S\relax
\newcommand{\S}{\mathbb{S}}
\DeclareMathOperator{\dist}{dist}
\DeclareMathOperator{\dom}{dom}
\DeclareMathOperator{\diag}{diag}
\DeclareMathOperator{\ones}{\mathbbm{1}}
\newcommand{\minimizeOver}[3]{\begin{array}{rl}\underset{#1}{\ptxt{minimize}} & #2\\ \ptxt{subject to} & #3\end{array}}
\newcommand{\maximizeOver}[3]{\begin{array}{rl}\underset{#1}{\ptxt{maximize}} & #2\\ \ptxt{subject to} & #3\end{array}}
\newcommand{\minimizationProblem}[2]{\minimizeOver{}{#1}{#2}}
\newcommand{\maximizationProblem}[2]{\maximizeOver{}{#1}{#2}}
\newcommand{\minimizeOverUnconstrained}[2]{\begin{array}{rl}\underset{#1}{\ptxt{minimize}} & #2\end{array}}
\newcommand{\maximizeOverUnconstrained}[2]{\begin{array}{rl}\underset{#1}{\ptxt{maximize}} & #2\end{array}}
\newcommand{\minimizationUnconstrained}[1]{\minimizeOverUnconstrained{}{#1}}
\newcommand{\maximizationUnconstrained}[1]{\maximizeOverUnconstrained{}{#1}}
\DeclareMathOperator{\argmin}{argmin}
\DeclareMathOperator{\argmax}{argmax}

% Proofs
\newcommand{\st}{s.t.}
\newcommand{\unique}{!}
\newcommand{\iffdef}{\overset{\ptxt{def}}{\Leftrightarrow}}
\newcommand{\eqVertical}{\rotatebox[origin=c]{90}{=}}
\newcommand{\mapsfrom}{\mathrel{\reflectbox{\ensuremath{\mapsto}}}}
\newcommand{\mapsdown}{\rotatebox[origin=c]{-90}{$\mapsto$}\mkern2mu}
\newcommand{\mapsup}{\rotatebox[origin=c]{90}{$\mapsto$}\mkern2mu}
\newcommand{\from}{\!\mathrel{\reflectbox{\ensuremath{\to}}}}
\newcommand{\labeledeq}[1]{\overset{\mathclap{\ptxt{#1}}}{=}}
\newcommand{\eqdef}{\labeledeq{def}}

% Brackets
\newcommand{\paren}[1]{\left(#1\right)}
\renewcommand{\brack}[1]{\left[#1\right]}
\renewcommand{\brace}[1]{\left\{#1\right\}}
\newcommand{\ang}[1]{\left<#1\right>}

% Algorithms
\algrenewcommand{\algorithmiccomment}[1]{\hskip 1em \texttt{// #1}}
\algrenewcommand\algorithmicrequire{\textbf{Input:}}
\algrenewcommand\algorithmicensure{\textbf{Output:}}
\newcommand{\algP}{\ptxt{\textbf{P}}}
\newcommand{\algNP}{\ptxt{\textbf{NP}}}
\newcommand{\algNPC}{\ptxt{\textbf{NP-Complete}}}
\newcommand{\algNPH}{\ptxt{\textbf{NP-Hard}}}
\newcommand{\algEXP}{\ptxt{\textbf{EXP}}}
\DeclareMathOperator{\fl}{fl}

%%%%%%%%%%%%%%%%%%%%%%%%%%%%%
% Other commands and shorthand
%%%%%%%%%%%%%%%%%%%%%%%%%%%%%
\newcommand{\flag}[1]{\textbf{\textcolor{red}{#1}}}
\let\uSym\u
\let\u\relax
\newcommand{\u}[1]{\underline{#1}}
\let\bSym\b
\let\b\relax
\newcommand{\b}[1]{\textbf{#1}}
\let\iSym\i
\let\i\relax
\newcommand{\i}[1]{\textit{#1}}
\let\scSym\sc
\let\sc\relax
\newcommand{\sc}[1]{\textsc{#1}}
\newcommand{\mf}[1]{\mathfrak{#1}}

%%%%%%%%%%%%%%%%%%%%%%%%%%%%%%%%%%%%%%%
% Make l's curvy in math environments %
%%%%%%%%%%%%%%%%%%%%%%%%%%%%%%%%%%%%%%%
\mathcode`l="8000
\begingroup
\makeatletter
\lccode`\~=`\l
\DeclareMathSymbol{\lsb@l}{\mathalpha}{letters}{`l}
\lowercase{\gdef~{\ifnum\the\mathgroup=\m@ne \ell \else \lsb@l \fi}}%
\endgroup

%%%%%%%%%%%%%%%%%%%%%%%%%
% Fix \vdots and \ddots %
%%%%%%%%%%%%%%%%%%%%%%%%%
\usepackage{letltxmacro}
\LetLtxMacro\orgvdots\vdots
\LetLtxMacro\orgddots\ddots

\makeatletter
\DeclareRobustCommand\vdots{%
	\mathpalette\@vdots{}%
}
\newcommand*{\@vdots}[2]{%
	% #1: math style
	% #2: unused
	\sbox0{$#1\cdotp\cdotp\cdotp\m@th$}%
	\sbox2{$#1.\m@th$}%
	\vbox{%
		\dimen@=\wd0 %
		\advance\dimen@ -3\ht2 %
		\kern.5\dimen@
		% remove side bearings
		\dimen@=\wd2 %
		\advance\dimen@ -\ht2 %
		\dimen2=\wd0 %
		\advance\dimen2 -\dimen@
		\vbox to \dimen2{%
			\offinterlineskip
			\copy2 \vfill\copy2 \vfill\copy2 %
		}%
	}%
}
\DeclareRobustCommand\ddots{%
	\mathinner{%
		\mathpalette\@ddots{}%
		\mkern\thinmuskip
	}%
}
\newcommand*{\@ddots}[2]{%
	% #1: math style
	% #2: unused
	\sbox0{$#1\cdotp\cdotp\cdotp\m@th$}%
	\sbox2{$#1.\m@th$}%
	\vbox{%
		\dimen@=\wd0 %
		\advance\dimen@ -3\ht2 %
		\kern.5\dimen@
		% remove side bearings
		\dimen@=\wd2 %
		\advance\dimen@ -\ht2 %
		\dimen2=\wd0 %
		\advance\dimen2 -\dimen@
		\vbox to \dimen2{%
			\offinterlineskip
			\hbox{$#1\mathpunct{.}\m@th$}%
			\vfill
			\hbox{$#1\mathpunct{\kern\wd2}\mathpunct{.}\m@th$}%
			\vfill
			\hbox{$#1\mathpunct{\kern\wd2}\mathpunct{\kern\wd2}\mathpunct{.}\m@th$}%
		}%
	}%
}
\makeatother

\newcommand{\B}{
	\begin{tikzpicture}
	\filldraw [fill=red, draw=black] (0, 0) rectangle (0.37, 0.45);
	\draw [line width=0.5mm, white ] (0.1,0.08) -- (0.1,0.38);
	\draw[line width=0.5mm, white ] (0.1, 0.35) .. controls (0.2, 0.35) and (0.4, 0.2625) .. (0.1, 0.225);
	\draw[line width=0.5mm, white ] (0.1, 0.225) .. controls (0.2, 0.225) and (0.4, 0.1625) .. (0.1, 0.1);
	\end{tikzpicture}
}

% Allow custom symbols as arrows in tikz-cd
\tikzset{
  symbol/.style={
    draw=none,
    every to/.append style={
      edge node={node [sloped, allow upside down, auto=false]{$#1$}}}
  }
}
% Example: \arrow[r,symbol=\cong]
% https://tex.stackexchange.com/questions/394154/how-to-include-inclusion-subgroup-relationship-in-tikz-cd-diagram

\author{Professor Alejandro Uribe-Ahumada\\ \small\i{Transcribed by Thomas Cohn}}
\title{Math 635 Lecture 9}
\date{2/8/21} % Can also use \today

\begin{document}
\maketitle
\setlength\RaggedRightParindent{\parindent}
\RaggedRight

\subsection*{An Example of Connections and Parallelism}

\par\noindent
Take $M=S^{1}$, with the coordinate $0\le{}x\le{}2\pi$. Let $\mathcal{E}=M\times\R^{2}=S^{1}\times\R^{2}$, with the global frame $(E_{1}=(1,0),E_{2}=(0,1))$, $\forall{}x\in{}S^{1}$. Then elements of $\Gamma(\mathcal{E})$ are functions $\vec{f}(x)=(f^{1}(x),f^{2}(x))$.\n

\par\noindent
Suppose $\nabla$ is given by $\vartheta$, any $2\times{}2$ matrix of $1$-forms on $M$. So $\theta_{i}^{j}=\alpha_{i}^{j}dx$, for $\alpha_{i}^{j}\in{}C^{\infty}(S^{1})$. Thus,
\[
	\nabla_{\pd{}{x}}\vec{f}=\underbrace{\begin{pmatrix}
			\frac{df^{1}}{dx}\\[5pt]
			\frac{df^{2}}{dx}
		\end{pmatrix}}_{=\frac{d}{dx}\vec{f}}+\underbrace{\begin{pmatrix}
		a_{i}^{j}
	\end{pmatrix}}_{\hspace{-20pt}\mathrlap{\ptxt{because $dx\paren{\pd{}{x}}=1$}}}\vec{f}
\]

\par\noindent
Consider the simple case $\vartheta=\begin{pmatrix}
	0 & c\\ -c & 0
\end{pmatrix}$, for some $c\in\R\cut\set{0}$. Then
\[
	\vec{f}\;\ptxt{is parallel}\quad\Leftrightarrow\quad\nabla_{\pd{}{x}}\vec{f}=0\quad\Leftrightarrow\quad
	\left\{\begin{array}{l}
		\frac{df^{1}}{dx}=cf^{2}\\[5pt]
		\frac{df^{2}}{dx}=-cf^{1}
	\end{array}\right.\quad\Leftrightarrow\quad\left\{\begin{array}{l}
		f^{1}(x)=f^{1}(0)\cos(cx)+f^{2}(0)\sin(cx)\\[5pt]
		f^{2}(x)=-f^{1}(0)\sin(cx)+f^{2}(0)\cos(cx)
	\end{array}\right.
\]
This final object is the solution to the parallel transport problem, starting at $x=0$ with $\smallPMatrix{f^{1}(0)\\ f^{2}(0)}=\vec{f}(0)$.
\[
	\begin{tikzpicture}[scale=1]
		\draw (0,0) circle [x radius=5cm, y radius=3cm];
		\draw (-4.5,0) node {$S^{1}$};
		%
		\node (A) at (4,-{sqrt(9*(1-(16/25)))}) {};
		\node (B) at (3,{sqrt(9*(1-(9/25)))}) {};
		%
		\fill (A) circle [radius=0.1cm];
		\draw (A) + (-1,1.2) -- +(-1,-0.8) -- +(1,-1.2) -- +(1,0.8) -- cycle;
		\draw[blue] [->] (A) edge +(0.75,{0.75*(-1.2 - (-0.8))/2});
		\draw[blue] [->] (A) edge +(0,0.75);
		\draw[red] [->] (A) edge +(0.9, 0.5);
		\draw (A) + (0,-0.5) node {$x=0$};
		\draw[red] (A) + ({0.9+0.5}, 0.5) node {$\vec{f}(0)$};
		\draw (A) + (1.25,-1.45) node {$\R^{2}$};
		%
		\fill (B) circle [radius=0.1cm];
		\draw (B) + (-1,-1.2) -- +(-1,0.8) -- +(+1,1.2) -- +(1,-0.8) -- cycle;
		% \draw[blue] [->] (B) edge +(0,0.75);
		% \draw[blue] [->] (B) edge +(-0.75,{-0.75*(-0.8 - (-1.2))/2});
		\draw[red] [->] (B) edge +(-0.5, 0.9);
		\draw (B) + (0,-0.5) node {$x=x_{1}$};
		\draw[red] (B) + (-0.5,{0.9+0.5}) node {$\vec{f}(x_{1})$};
		\draw (B) + (1.25,1.45) node {$\R^{2}$};
	\end{tikzpicture}
\]
Additionally, we can map from $\vec{f}(0)$ to $\vec{f}(2\pi)$ with the holonomy matrix:
\[
	\begin{pmatrix}
		\cos(2\pi{}c) & -\sin(2\pi{}c)\\
		\sin(2\pi{}c) & \cos(2\pi{}c)
	\end{pmatrix}
\]
Thus, the holonomy is trivial (i.e., the holonomy matrix is the identity matrix $I_{2}$) iff $c\in\Z$.\n

\exer{
	Compute $\nabla_{\pd{}{x}}E_{1}$ and $\nabla_{\pd{}{x}}E_{2}$. Note: It won't be $0$!\n
}

\subsection*{Curvature of a Connection on a Bundle}

\par\noindent
Say $\nabla$ is a connection on the vector bundle $\mathcal{E}\to{}M$. Curvature is a local object, and there are two approaches to describe it.\n

\subsubsection*{First Approach}

\par\noindent
Begin with a moving frame $(E_{1},\ldots,E_{r})$ (which leads to the connection matrix $\vartheta=(\theta_{i}^{j})$).\n

\par\noindent
Question: Can we change the moving frame to $(F_{1},\ldots,F_{r})$ so that the new connection matrix, $\tilde\vartheta$, is $0$?\n

\par\noindent
Well, if $A=(a_{j}^{i})$ is defined so that $F_{j}=a_{j}^{i}E_{i}$ (using Einstein summation notation), then we saw (in homework) that
\[
	\tilde\vartheta=A\inv{}dA+A\inv\vartheta{}A
\]
So
\[
	\tilde\vartheta=0\quad\Leftrightarrow\quad{}dA=-\vartheta{}A\quad\Leftrightarrow\quad\forall{}i,j,da_{j}^{i}=-\theta_{k}^{i}a_{j}^{k}
\]
We suspect that we'll encounter problems from the fact that $d^{2}=0$. Assume such a solution $A$ exists. Then, using the product rule, $0=-(d\theta_{k}^{i})a_{j}^{k}+\theta_{k}^{i}\wedge{}da_{j}^{k}$. But we know $da_{j}^{k}=-\theta_{l}^{k}a_{j}^{l}$, so for the connection matrix to be trivial, we must have
\[
	(d\theta_{k}^{i})a_{j}^{k}+\theta_{k}^{i}\wedge\theta_{l}^{k}a_{j}^{l}=0,\quad\forall{}i,j
\]
Note: $\theta_{k}^{j}\wedge\theta_{l}^{k}=(\vartheta\wedge\vartheta)_{l}^{j}$ by definition. So using matrix notation, we can write this necessary condition as
\[
	(d\vartheta)A+(\vartheta\wedge\vartheta)A=0
\]
$A$ is invertible, so this condition is true iff $d\vartheta+\vartheta\wedge\vartheta=0$.\n

\defn{
	$\Omega\eqdef{}d\vartheta+\vartheta\wedge\vartheta$ is a matrix of $2$-forms, called the \u{curvature matrix} of $\nabla$ with respect to the moving frame $(E_{1},\ldots,E_{r})$.
}

\par\noindent
The connection is trivializable only if $\Omega=0$ (i.e. this is a necessary condition). There is no reason in principle this should happen, so it's very rare!

\subsubsection*{Second Approach}

\par\noindent
Fix $X,Y\in\mathfrak{X}(M)$. Consider $\nabla_{X},\nabla_{Y}:\Gamma(\mathcal{E})\to\Gamma(\mathcal{E})$ as operators, acting linearly on smooth sections.\n

\par\noindent
Question: What is their commutator (as operators)?\n

\par\noindent
Well, if $\mathcal{E}=M\times\R^{r}$ and $\vartheta\equiv{}0$, then $\nabla_{X}\vec{f}=X(f)+0$. So $\nabla_{X}=X$. Therefore, $\brack{\nabla_{X},\nabla_{Y}}=\brack{X,Y}=\nabla_{\brack{X,Y}}$, by definition of the commutator of vector fields.\n

\par\noindent
In general, let $E_{1},\ldots,E_{r}$ be a moving frame, and $\vartheta=(\theta_{i}^{j})$ the connection matrix as before. Then $\nabla_{X}=X+\vartheta(X)$, acting on functions $\vec{f}:U\to\R^{r}$ ($\vartheta(X)$ is a matrix of functions). Likewise, $\nabla_{Y}=Y+\vartheta(Y)$. So we can compute:
\[
	\brack{\nabla_{X},\nabla_{Y}}=\brack{X+\vartheta(X),Y+\vartheta(Y)}=\brack{X,Y}+\brack{\vartheta(X),\vartheta(Y)}+\underbrace{\brack{X,\vartheta(Y)}}_{\ptxt{(I)}}+\underbrace{\brack{\vartheta(X),Y}}_{\ptxt{(II)}}
\]
For (I), we have
\[
	\paren{\brack{X,\vartheta(Y)}(\vec{f})}^{i}=X\underbrace{(\theta_{j}^{i}(Y)f^{j})}_{\mathclap{\substack{
		\ptxt{$i$th component}\\
		\ptxt{of $\vartheta(Y)(\vec{f})$}
	}}}-\theta_{j}^{i}(Y)X(f^{j})=(X\theta_{j}^{i}(Y))f^{j}
\]
by Leibniz's rule. So (I) is $X((\theta_{j}^{i})(Y))=X(\vartheta(Y))$ as a matrix. Similarly, (II) is $-Y((\theta_{j}^{i})(X))=-Y(\vartheta(X))$ as matrix.\n

\par\noindent
Claim: This eventually leads to $\brack{\nabla_{X},\nabla_{Y}}=\nabla_{\brack{X,Y}}+d\vartheta(X,Y)+\brack{\vartheta(X),\vartheta(Y)}$.\n
Proof: Use the fact that for any $1$-form $\omega$, and for any $X,Y\in\mf{X}(M)$, $d\omega(X,Y)=X(\omega(Y))-Y(\omega(X))-\omega(\brack{X,Y})$.\n

\par\noindent
Observe: $\brack{\vartheta(X),\vartheta(Y)}=(\vartheta\wedge\vartheta)(X,Y)$, by definition of $\wedge$. So we get the same object using both approaches!\n

\defn{
	$\mathcal{R}(X,Y)\eqdef{}d\vartheta(X,Y)+\brack{\vartheta(X),\vartheta(Y)}$ is the \u{curvature operator}.\n
}

\end{document}